% ──────────────────────────────────────────────────────────────────────────────
%  Section 2 — The sending current as an affine functional of S
% ──────────────────────────────────────────────────────────────────────────────

\section{\texorpdfstring{The phase current as an affine functional of $S$}{The phase current as an affine functional of S}}
\label{sec:affine}

Substituting \autoref{eq:W-recall} into the sending-current relation
$\vI_{s} = \matB^{-1}(W^{*}\vu - V_{s}\,\matA\vu)$ and reading off the
$k$-th component yields, after a straightforward algebraic
rearrangement,
%
\begin{equation}
\label{eq:Ik-affine}
I_{k}(S,\ell) \;=\; a_{k}(\ell)\,S^{*} \;+\; b_{k}(\ell),
\qquad k \in \{A,B,C\},
\end{equation}
%
where the two complex coefficients depend on the line only:
%
\begin{equation}
\label{eq:ak-bk}
a_{k}(\ell) \;=\;
   \frac{p_{k}(\ell)}{V_{s}\,\overline{\alpha(\ell)}},
\qquad
b_{k}(\ell) \;=\;
   \frac{V_{s}\,\bigl[\,p_{k}(\ell)\,\overline{\beta(\ell)} \;-\;
                          \overline{\alpha(\ell)}\,q_{k}(\ell)\,\bigr]}
        {\overline{\alpha(\ell)}}.
\end{equation}

\noindent Each phase current is therefore an affine functional of the
complex apparent power: $a_{k}(\ell)\in\mathbb{C}$ is the gain on the
loading term $S^{*}$, and $b_{k}(\ell)\in\mathbb{C}$ is the no-load
offset that the line delivers at $S=0$. Both coefficients are produced
once per length from \autoref{eq:line-constants} and
\autoref{eq:alpha-beta}, and the dependence on $S$ is fully linear.

The affine form \autoref{eq:Ik-affine} is exact within the model
assumptions of the parent document (symmetric local-bus voltage,
balanced remote-bus magnitude, lumped end-to-end $\matM(\ell)$) and
absorbs every detail of the non-circulant per-unit-length matrices
into the six complex line constants
$\{a_{A},a_{B},a_{C},b_{A},b_{B},b_{C}\}(\ell)$.
